% Original training problem based on deidentified topic suggestions.
\documentclass[11pt,a4paper]{article}
\usepackage[margin=23mm,headheight=14pt]{geometry}
\usepackage{fontspec}
\setmainfont{Noto Serif CJK SC}[ItalicFont={Noto Serif CJK SC},ItalicFeatures={FakeSlant=0.15}]
\setsansfont{Noto Sans CJK SC}\setmonofont{DejaVu Sans Mono}[Scale=.81]
\XeTeXlinebreaklocale "zh"\XeTeXlinebreakskip=0pt plus 1pt
\usepackage{amsmath,amssymb,booktabs,array,enumitem,tikz,fancyhdr,hyperref,xurl}
\hypersetup{unicode=true,colorlinks=true,linkcolor=black,urlcolor=blue!50!black,pdftitle={校园小巴的高峰运行与调度},pdfauthor={SZUMMA2026 Training}}
\setlength{\parindent}{2em}\setlength{\parskip}{5pt}
\setlist[enumerate]{leftmargin=2em,itemsep=10pt,topsep=7pt}
\pagestyle{fancy}\fancyhf{}\fancyhead[L]{\small 数学建模训练题\quad Problem\_005}\fancyhead[R]{\small 审核稿}\fancyfoot[C]{\thepage}
\renewcommand{\headrulewidth}{.3pt}\newcommand{\file}[1]{\texttt{\detokenize{#1}}}
\begin{document}
\begin{center}{\LARGE\bfseries 校园小巴的高峰运行与调度}\end{center}

校园小巴连接教学区、宿舍区、食堂和校门，为师生提供短途出行服务。上下课前后，乘车需求往往集中出现，站点之间的客流也会随时段变化。一辆车在某站接走的乘客越多，停站时间通常越长；乘客在何处下车，又决定了后续各站可用的座位。发车安排、乘降和车辆周转因而相互影响。

现结合深圳大学粤海、丽湖两校区的功能分布，研究两条校内环线的高峰运行安排。两校区各自配置车队，独立运行。希望在给定车辆和运行时间内，提高乘客按时到达的比例，改善候车较长站点的服务。

粤海环线记为Y，设8个站点、5辆小巴；丽湖环线记为L，设9个站点、4辆小巴。站点从0编号，车辆按 $0\to1\to\cdots\to n-1\to0$ 的固定站序循环，每站必停。图中给出了站点连接关系，各站对应的功能区和相邻站间行时见附件。

\begin{center}
\begin{minipage}{.48\linewidth}\centering
\begin{tikzpicture}[>=stealth,font=\small]
\foreach \i in {0,...,7}{\pgfmathsetmacro{\a}{90-45*\i}\coordinate (y\i) at (\a:1.65);}
\foreach \i in {0,...,7}{\pgfmathtruncatemacro{\j}{mod(\i+1,8)}\draw[->,shorten <=11pt,shorten >=11pt] (y\i)--(y\j);}
\foreach \i in {0,...,7}{\node[circle,draw,fill=white,inner sep=3pt] at (y\i) {Y\i};}
\foreach \i in {0,4}{\node[circle,draw,fill=orange!20,inner sep=3pt] at (y\i) {Y\i};}
\node at (0,-2.2) {粤海环线Y};
\end{tikzpicture}
\end{minipage}\hfill
\begin{minipage}{.48\linewidth}\centering
\begin{tikzpicture}[>=stealth,font=\small]
\foreach \i in {0,...,8}{\pgfmathsetmacro{\a}{90-40*\i}\coordinate (l\i) at (\a:1.65);}
\foreach \i in {0,...,8}{\pgfmathtruncatemacro{\j}{mod(\i+1,9)}\draw[->,shorten <=11pt,shorten >=11pt] (l\i)--(l\j);}
\foreach \i in {0,...,8}{\node[circle,draw,fill=white,inner sep=3pt] at (l\i) {L\i};}
\foreach \i in {0,4}{\node[circle,draw,fill=orange!20,inner sep=3pt] at (l\i) {L\i};}
\node at (0,-2.2) {丽湖环线L};
\end{tikzpicture}
\end{minipage}

\small 站序示意图。着色站点0、4为控制站；Y0表示粤海环线的0站，L0同理。
\end{center}

每车设14个乘客座位，载客量以14人为上限。初始各站候车人数为零，所有车辆空车停在本校区0站的场站内。研究中忽略道路和站台阻塞，车辆可以交换先后次序，控制站有足够停留空间。每条出行请求对应一名乘客，其到站时刻、起点、目的站和到达期限见附件。乘客在给定起点持续候车，乘一辆车到给定目的站完成行程。

车辆到站先下后上，本次上车乘客从到站瞬间的队列中按先到先上选取，选中后转入上车服务；随后到站的乘客等候下一车。若本次下车、上车人数分别为 $A、B$，站停时间为
\[
\tau=10+2A+3B\quad\text{秒}.
\]
下车者统一在到站后 $10+2A$ 秒完成下车，上车者统一在站停结束时完成上车；候车时间计到上车完成。同秒旅客先入队，再处理车辆到站；同秒多车按车辆编号先后分配乘客，详细事件约定见附件。

\newpage
\begin{center}{\large\bfseries 运行安排与服务评价}\end{center}

\textbf{运行计划。} 各车首次驶入0站开始服务的时刻可在第0至600秒间选择，完成乘降后离站。此后每次回到0站，可在关门后按计划固定等候0至60秒，首次除外。对照方案为同一校区第 $k$ 辆车在 $120(k-1)$ 秒首次投入，各站有位即上，关门后随即离站，完成一圈后继续运行。

\textbf{运行调整。} 为改善后续站点的服务，可在车辆到站时设置0至4个预留座位。在本次下车后，设仍在车人数为 $L$，本站尚待分配的候车人数为 $Q$，预留数为 $r$，则本次上车人数为
\[
B=\min\{Q,\max(14-L-r,0)\}.
\]
乘客仍按队列顺序上车，在车乘客均在各自目的站下车。在0、4站，还可安排关门后附加等待0至90秒，称为控车等待；等候期间车门关闭。动态控车取代原计划中的固定等待，其余站点在乘降结束后离站。时间以秒计，投入时刻、等待时长和预留座位数均取整数。

路段行时为附件基础行时、当前司机系数和进入路段时的工况系数之积，向上取整，不含站停。车辆进入路段时确定本段行时。各司机的行时系数在给定工作阶段内保持为附件值。

\textbf{调度信息。} 每校区提供P1、P2、P3三份规划样本，用于研究需求规律和制定运行计划。实时调度有 $\delta$ 秒信息延迟，基准为 $\delta=60$。在时刻 $t$，本车身份、所在站、本次下车人数、下车后余载和本站当前队列人数即时可得；其余信息来自时间戳不晚于 $t-\delta$ 的运行报告。报告记录车辆到离站时刻、到站当时的队列人数、上下客数和总载客量，按事件实际发生时刻形成。延迟为0时，同秒报告按事件处理顺序依次可用。

初态、基础行时、司机系数、乘降参数及交班窗口事前给定。控制器的输入由这些条件、上述即时量和延迟报告组成，可利用规划样本预测需求。附件中的逐人起讫站、期限、完整工况记录和工况编号用于运行回放与事后评价。同一校区的N、R、E三个检验工况在第1800秒以前具有相同的外部输入，相应决策应一致。

\textbf{交班。} 第四问考虑司机交接。每辆车配有独立接班人员，在给定窗口内于0站待命。车辆到站时一并决定预留、控车及本次是否交班。交班在该车于0站完成本次乘降后立即开始，持续120秒，整个过程位于给定窗口内，车内乘客留车。与控车等待重叠时，额外停留取两者较大值，完成后改用接班司机系数。每车在窗口内交班一次，窗口结束后由接班司机继续运行。

\textbf{服务评价。} 从第0秒观察至 $T=7200$ 秒，乘客在前5400秒到站。到达期限已预留后续步行时间，以完成下车为准。按时到达率为观察期内按期完成下车的人数占全部请求的比例，另统计观察结束时的未完成人数。设 $t_{\mathrm{arrival}}$ 为到站时刻，$t_{\mathrm{board}}$ 为上车完成时刻，每名乘客的候车耗时为
\[
W=\min\{t_{\mathrm{board}},T\}-t_{\mathrm{arrival}},
\]
未完成上车时以 $T$ 代入。报告全体平均值、各起点的95分位及最不利起点；95分位取该起点 $m$ 个样本升序后的第 $\lceil0.95m\rceil$ 个数。对已上车者，乘车耗时从完成上车计至完成下车，未完成下车者计至 $T$，同时报告其人数。

\newpage
\begin{center}{\large\bfseries 研究任务}\end{center}

请建立数学模型，研究以下问题。
\begin{enumerate}
\item 分别建立两条环线车辆与乘客的运行模型，分析对照方案的服务情况。从运力总量、目的站分布和乘降停站等方面解释候车与到达结果，区分运力不足和运行安排的影响。选择一种简化模型作比较，讨论其误差与适用范围，并检验人数守恒及座位容量。
\item 保持各校区车队、线路和运行时长，利用P1、P2、P3制定一套事前运行计划，确定各车首次投入时刻和回到0站的固定等待。本问有位即上，附加等待按事前计划执行。将计划用于N、R、E工况，与对照方案比较，并说明两校区安排的异同。
\item 在第二问计划的基础上，提出信息延迟下的控车与座位预留规则，比较按时到达率、候车、在车耗时和各起点的服务。研究突发客流和行时延长的影响，比较信息延迟为0、60、120秒时的结果，并检验相同输入历史下的决策一致性。说明各类调整的作用及其适用条件。
\item 加入司机交班要求，研究第二、三问方案的可行性，作出必要调整，并分析交班对服务的影响。为两校区分别拟订小规模试运行计划，确定需要校准的参数，说明数据采集与观测安排，提出检验模型和评价试运行效果的方法。
\end{enumerate}

请在论文中说明模型假设、运行安排及其依据，报告各项服务指标，并附可复现程序。可根据所分析的服务目标确定综合评价方法，给出面向校园交通管理部门的运行建议。

\bigskip
\noindent\textbf{工况说明}

P1、P2、P3为规划样本。N为正常检验工况，R、E由N分别增加行时和客流得到：R从第1800秒起将路段系数乘以1.25；E在第2100至2400秒于指定站点增加一批请求，其余输入与N相同。每个校区采用同一套运行计划和调度规则检验三个工况。可另行生成样本考察方案的稳定性，生成条件在论文中说明。

\newpage
\begin{center}{\large\bfseries 参数与附件}\end{center}

\begin{center}
\begin{tabular}{lrr}
\toprule
项目 & 粤海环线Y & 丽湖环线L\\\midrule
站点数 & 8 & 9\\
小巴数 & 5 & 4\\
每车乘客座位 & 14 & 14\\
规划样本数 & 3 & 3\\
检验工况数 & 3 & 3\\
观察时长（秒） & 7200 & 7200\\\bottomrule
\end{tabular}
\end{center}

\begin{center}
\begin{tabular}{p{.39\linewidth}p{.54\linewidth}}
\toprule
文件 & 内容\\\midrule
\file{stops.csv}、\file{links.csv} & 功能节点、允许连接及基础路段行时\\
\file{vehicles.csv} & 车队、容量、交班前后司机系数及窗口\\
\file{cases.csv} & 规划样本、检验工况及其配对关系\\
\file{passenger_requests.csv} & 乘客到站时刻、起讫站和完成下车的期限\\
\file{link_factors.csv} & 各工况分时段的路段行时系数\\
\file{baseline_plan.csv} & 两校区对照投入安排\\
\file{rules.json}、\file{data/README.md} & 运行规则、同刻事件顺序、字段单位、信息及评价口径\\\bottomrule
\end{tabular}
\end{center}

附件含12个运行样本、7350条出行请求及2448条路段系数记录。路段行时、乘降及等待共同决定车辆的到离站时刻；更换运行安排后，应按相应乘降人数重新计算站停。字段、单位和事件顺序见\file{data/README.md}。

本题为原创模拟题，线路、参数和工况采用模拟设定。数据生成方法、来源说明及补测字段见附件说明。

\noindent\textbf{参考资料}

\small
深圳大学\href{https://www.szu.edu.cn/xxgk/xydt.htm}{校园地图}。

公交运行间隔与控制的研究可参见 Daganzo，2009，\href{https://doi.org/10.1016/j.trb.2009.04.002}{A headway-based approach to eliminate bus bunching}；Xuan、Argote、Daganzo，2011，\href{https://doi.org/10.1016/j.trb.2011.07.009}{Dynamic bus holding strategies for schedule reliability}。

建议比赛时间为96小时。
\end{document}
